\documentclass[UTF8,a4paper,12pt]{ctexart}

\usepackage{amsmath,amssymb,bm}
\usepackage{booktabs}
\usepackage{geometry}
\usepackage{longtable}
\usepackage{array}
\usepackage{enumitem}
\usepackage{microtype}
\usepackage{hyperref}

\geometry{left=2.6cm,right=2.6cm,top=2.5cm,bottom=2.5cm}
\setlength{\parindent}{2em}
\setlength{\parskip}{0.35em}
\renewcommand{\arraystretch}{1.35}
\setlist{nosep}
\hypersetup{hidelinks}

\title{方法与案例研究部分}
\author{}
\date{}

\begin{document}

\maketitle

本文方法与案例研究部分围绕一条主线展开：大型多线换乘站应急疏散中，多线路出站客流、列车释放客流、站厅滞留客流和换乘客流会同时进入站内网络，并在楼扶梯、换乘通道、闸机和出口前区等共享设施处形成排队瓶颈；因此，疏散组织不能只寻找几何上的最短可达路径，还需要在路径选择时考虑乘客到达共享设施时的排队等待。基于这一问题，本文先构建站内网络流与共享设施容量约束模型，再利用 AA 方法进行到达时刻队列感知的路径组织，最后通过低负荷、高负荷、Pathfinder 对照和敏感性分析验证全站疏散效率、瓶颈缓解和适用边界。

上述逻辑可概括为：多线换乘站多源客流在共享设施处形成疏散瓶颈；构建站内网络流和共享容量约束；预测客流批次到达共享设施时的排队等待；用 AA 求解队列感知路径组织方案；用仿真场景验证疏散效率、设施瓶颈和方法边界。第 3 章回答“疏散问题如何建模”，第 4 章回答“AA 如何在该模型上求解路径”，第 5 章回答“这种路径组织是否改善了全站疏散”。

\section{3 基于站内网络流的换乘站疏散模型}

龙阳路站这类大型多线换乘站的应急疏散，本质上可以看作一个带共享容量约束的站内网络流问题。乘客从站台、列车、站厅和换乘入口等不同来源进入网络，沿通道、楼扶梯、闸机和出口向安全区域移动。问题的关键不是网络是否连通，而是多来源客流在有限设施上的叠加会改变实际通行时间。若路径组织只依据当前可见阻抗或几何距离，部分客流可能在真正到达设施时遇到由前序客流形成的队列，从而导致尾部清空时间延长。

因此，本章只做一件事：把车站疏散写成“客流如何在站内网络中流动，并受哪些容量约束控制”。AA 不是本章的模型对象，而是第 4 章用于求解该网络流路径组织问题的方法。

\subsection{3.1 站内疏散网络与多源客流输入}

将车站公共区抽象为有向网络 \(G=(V,E)\)。节点 \(V\) 表示站台等待区、列车释放点、站厅区域、换乘通道节点、闸机区、出口前区和安全出口；边 \(E\) 表示乘客可通行的连接关系。对于楼梯、扶梯、换乘通道、闸机和出口等具有实际通行能力的设施，设其资源集合为 \(R\)。若多条路径使用同一物理设施，则通过映射
\[
\rho:E\rightarrow R\cup\varnothing
\]
将相关边归并到同一资源 \(r\in R\)，使其共享同一服务能力。

客流输入按来源组生成。设来源组集合为
\[
\mathcal G=\mathcal G^{p}\cup \mathcal G^{train}\cup \mathcal G^{hall}\cup \mathcal G^{tr},
\]
其中，\(\mathcal G^{p}\) 表示站台等待客流，\(\mathcal G^{train}\) 表示列车到达释放客流，\(\mathcal G^{hall}\) 表示站厅滞留客流，\(\mathcal G^{tr}\) 表示换乘客流。换乘客流不是本文单独优化的对象，但它是大型换乘站整体疏散中的关键组成部分：它通常跨越不同线路空间，经过的共享设施更多，更容易与其他来源客流在同一楼扶梯、通道或闸机前汇合。因此，本文保留来源组标识，用于后续分析换乘相关客流如何参与全站瓶颈形成。

在离散时间 \(t\) 下，可参与路径组织的客流批次记为
\[
b=(g_b,o_b,D_b,n_b,t_b),
\]
其中，\(g_b\) 为来源组，\(o_b\) 为当前位置，\(D_b\) 为可选安全出口集合，\(n_b\) 为人数，\(t_b\) 为进入网络的时刻。批次化处理的意义在于保留列车集中到达、换乘释放和站厅滞留之间的时间差异，使模型能够分析“何时到达设施”对排队和疏散效率的影响。

\begin{table}[htbp]
\centering
\caption{站内网络流模型中的客流来源}
\small
\begin{tabular}{p{0.20\textwidth}p{0.30\textwidth}p{0.38\textwidth}}
\toprule
来源组 & 网络入口 & 疏散意义 \\
\midrule
站台等待客流 & 各线路站台等待区 & 形成基础出站需求，与列车释放流共同竞争垂向设施 \\
列车释放客流 & 各线路列车车厢或站台分区 & 在短时间内增加站台和楼扶梯负荷，是高负荷场景的主要来源 \\
站厅滞留客流 & 站厅 staging 节点 & 影响闸机、出口前区和站厅连接通道的接收状态 \\
换乘客流 & 换乘通道或跨线连接入口 & 跨越多线路空间，参与共享设施瓶颈和出口负荷分配 \\
\bottomrule
\end{tabular}
\end{table}

\subsection{3.2 共享设施容量与网络流约束}

设 \(F_e(t)\) 为时间步 \(t\) 内通过边 \(e\) 的客流量，\(Q_r(t)\) 为共享设施 \(r\) 入口处的等待队列，\(\mu_r\) 为设施服务率。时间步长为 \(\Delta t\) 时，设施 \(r\) 的可释放能力为
\[
C_r(t)=\left\lfloor \mu_r\Delta t+\xi_r(t)\right\rfloor,
\]
其中，\(\xi_r(t)\) 为离散化余量项。对所有映射到同一资源的边，当前时间步的通过量需要满足
\[
\sum_{e:\rho(e)=r}F_e(t)\leq C_r(t)+Q_r(t),\quad r\in R.
\]
这表示同一楼梯、扶梯、闸机或出口前区只能释放一次真实容量，不能因为网络中存在多条路径而重复计算通行能力。

资源队列按照到达和服务释放更新：
\[
Q_r(t+\Delta t)=
\max\left\{
Q_r(t)+A_r(t)-C_r(t),0
\right\},
\]
其中，\(A_r(t)\) 表示时间步 \(t\) 到达资源 \(r\) 的客流量。若多个来源组同时到达同一资源，它们共同进入 \(A_r(t)\)，并在同一服务能力下排队。

节点空间也需要满足接收能力。设 \(N_v(t)\) 为节点 \(v\) 的占用人数，\(\bar N_v\) 为最大接收人数，\(B_v(t)\) 为当前时间步预计进入节点 \(v\) 的人数，则有
\[
N_v(t)+B_v(t)\leq \bar N_v.
\]
当下游空间不能接收更多乘客时，上游乘客即使已经选择了路径，也不能继续被推进，而是在边末端、节点或设施入口形成等待。这样，模型既能表达设施服务能力不足造成的排队，也能表达下游空间拥挤造成的排队扩散。

全系统满足人数守恒：
\[
\sum_{v\in V}N_v(t)+\sum_{e\in E}M_e(t)+\sum_{r\in R}Q_r(t)+E^{out}(t)=N^{tot},
\]
其中，\(M_e(t)\) 为边上在途人数，\(E^{out}(t)\) 为累计离站人数，\(N^{tot}\) 为场景总人数。该约束保证所有乘客要么处于节点、边、队列中，要么已经完成疏散。

\subsection{3.3 疏散目标与评价指标}

在上述网络流模型中，路径组织的目标不是单纯缩短路径长度，而是在共享设施容量受限的条件下提高全站疏散效率并减少尾部滞留。对批次 \(b\) 的候选路径 \(p\)，其代价由行走时间、共享设施等待时间和拥挤暴露构成：
\[
C_{bp}(t)=T_{bp}^{walk}(t)+T_{bp}^{wait}(t)+\lambda R_{bp}(t).
\]
其中，\(T_{bp}^{walk}(t)\) 为路径行走时间，\(T_{bp}^{wait}(t)\) 为设施排队等待时间，\(R_{bp}(t)\) 为拥挤暴露项，\(\lambda\) 为尺度权重。第 4 章的 AA 方法重点改进 \(T_{bp}^{wait}(t)\) 的估计方式，即不只看当前队列，而是预测乘客到达设施时的队列。

本文从四类指标评价疏散效果。第一类是整体效率，包括 \(T_{50}\)、\(T_{80}\)、\(T_{95}\)、\(T_{99}\)、\(T_{100}\)、平均疏散时间和平均通过率，其中
\[
T_{\theta}=\min\{t\mid E^{out}(t)\geq \theta N^{tot}\}.
\]
第二类是瓶颈暴露，包括累计静止人秒、资源队列人秒、关键设施峰值队列和队列持续时间。第三类是负荷分配，包括出口通过量、关键设施通过量和 Jain 均衡指数。第四类是方法代价，包括总移动距离、平均移动时间和计算耗时。上述指标共同回答：路径组织是否缩短全站尾部清空，是否降低关键设施排队，是否只是把瓶颈转移到其他位置，以及是否付出了明显绕行或计算代价。

\section{4 基于 AA 的队列感知路径组织方法}

第 3 章把换乘站疏散表达为带共享容量约束的站内网络流问题。本章说明如何在该网络上组织乘客路径。本文采用两类方法进行比较：一类是当前状态 A* 对照方法，用于表示只依据当前可见网络阻抗进行路径搜索的计算基准；另一类是 AA 方法，它在 A* 搜索中加入到达时刻队列预测，使路径选择能够提前考虑前序客流对共享设施的占用。两种方法的差异只在路径组织层，执行层、网络、客流和容量约束保持一致。

\subsection{4.1 当前状态 A* 对照方法}

A* 是图搜索中的经典启发式最短路方法。对搜索节点 \(v\)，其优先级由
\[
f(v)=g(v)+h(v)
\]
确定。其中，\(g(v)\) 为从起点到当前节点的累计代价，\(h(v)\) 为从当前节点到目标出口的剩余代价估计。在本文中，当前状态 A* 对照方法使用当前节点占用、当前边权和当前可见通行阻抗评价路径，并在同一共享容量执行层下推进乘客。它的作用不是代表现场既定疏散方案，而是提供一个计算基准，用来检验“只看当前状态”与“预测到达时刻队列”之间的差异。

该对照方法的局限在于，它能反映当前是否拥堵，却难以反映乘客到达共享设施时的未来排队。例如，某一楼梯当前队列较短，但已有多个批次正在向该楼梯移动；当前状态 A* 仍可能继续把后续客流分配到该楼梯，导致执行层中出现排队叠加。本文设置 AA 方法，正是为了解决这一类由未来到达负荷引起的路径组织滞后问题。

\subsection{4.2 AA 到达时刻队列预测}

AA 的核心不是重新构造一个复杂模型，而是在 A* 搜索中改变路径代价的计算方式。对候选路径上的共享设施 \(r\)，AA 不只读取当前队列 \(Q_r(t)\)，还读取已经在途的客流和本轮已承诺的未来到达事件。设当前时刻为 \(t\)，资源 \(r\) 的未来到达事件集合为 \(\mathcal A_r(t)\)。若批次预计在 \(\theta\) 时刻到达该资源，则预测队列为
\[
\widehat Q_r(\theta|t)=
\max\left\{
Q_r(t)+
\sum_{(\eta_j,a_j)\in\mathcal A_r(t),\eta_j\leq \theta}a_j
-\mu_r(\theta-t),0
\right\}.
\]
对应的预测等待时间为
\[
\widehat W_r(\theta|t)=\frac{\widehat Q_r(\theta|t)}{\mu_r}.
\]
因此，一条几何上较短的路径如果在到达时刻面对较长队列，其综合代价可能高于另一条稍长但队列更短的路径。

AA 采用时间依赖多标签搜索。一个标签表示为
\[
\ell=(v,\eta,c,\pi),
\]
其中，\(v\) 为节点，\(\eta\) 为预计到达时刻，\(c\) 为累计代价，\(\pi\) 为前驱标签。同一节点在不同到达时刻可能面对不同的后续队列，因此 AA 不简单把所有到达同一节点的状态合并，而是保留必要的时间标签，并删除明显被支配的标签。沿边扩展时，若边对应共享设施，边代价中加入预测等待；若不对应共享设施，则主要加入行走时间。

\subsection{4.3 路径组织与共享容量执行}

在每个时间步，AA 对当前可调度批次逐一搜索路径。路径被接受后，系统登记该批次预计到达各共享设施的时间和人数，使后续批次能够看到前序批次对未来容量的占用。这样，AA 的自适应性不是简单地在拥堵后调高边权，而是在路径组织阶段提前计算“这批乘客到达设施时是否会排队”。

路径搜索完成后，所有路径意图都进入同一共享容量执行层。执行层按照设施服务能力、节点接收能力和在途移动过程释放乘客，并记录实际队列、停滞、出口通过量和来源组清空时间。若路径组织阶段预测准确，关键设施排队应在执行层中得到缓解；若预测不准确，下一时间步会根据更新后的状态重新搜索。由此，本文比较的是两种路径组织方法在同一物理约束下的执行效果，而不是比较两个脱离容量约束的最短路结果。

本文的求解流程可概括如下：
\begin{enumerate}
\item 读取当前节点占用、设施队列、在途客流和未来到达事件；
\item 生成当前时间步可调度客流批次；
\item 当前状态 A* 根据即时阻抗搜索路径，AA 根据到达时刻预测队列搜索路径；
\item 对 AA，记录被接受路径对共享设施的未来到达承诺；
\item 在共享容量执行层释放乘客并更新队列、节点占用和离站人数；
\item 重复上述过程，直至所有乘客完成疏散。
\end{enumerate}

\section{5 龙阳路站案例研究与仿真验证}

本章验证的不是“AA 这个算法名称是否更好”，而是第 3 章和第 4 章形成的完整机制是否成立：多来源客流是否会在共享设施处形成影响全站效率的排队瓶颈；AA 是否通过预测到达时刻队列改变路径分配；这种分配变化是否转化为更短的尾部清空、更低的停滞等待和更合理的出口利用；同时，方法是否带来绕行距离和计算时间增加。实验按这一证据顺序组织，而不是先堆总时间表。

\subsection{5.1 案例站与疏散场景}

龙阳路站包含多条轨道交通线路和磁浮线路，站内存在多处换乘通道、楼扶梯、闸机和出口前区。案例网络由 CAD 图纸和节点边关系建立。设施几何尺寸、服务能力和行人行为参数在正式定稿中应分为三类记录：CAD 可量测值、模型设定值和文献或规范取值。尚未核验来源的参数不写成实测值。

案例设置低负荷基础疏散和高负荷列车叠加疏散两个场景。低负荷场景包含站台等待、站厅滞留和换乘客流，用于检查共享设施竞争较弱时 AA 是否保持稳定；高负荷场景在基础客流上叠加多线路列车释放客流，用于检验多来源到达重叠下 AA 是否能够降低关键设施排队和尾部清空时间。根据当前模型配置，低负荷场景为 2187 人，高负荷场景为 17905 人；这些数值在结果冻结后应与最终输入文件再次核对。

\begin{table}[htbp]
\centering
\caption{案例场景与验证目的}
\small
\begin{tabular}{p{0.20\textwidth}p{0.33\textwidth}p{0.35\textwidth}}
\toprule
场景 & 客流构成 & 验证目的 \\
\midrule
低负荷基础疏散 & 站台等待、站厅滞留和换乘客流 & 检查瓶颈较弱时 AA 是否保持稳定，避免无必要绕行 \\
高负荷列车叠加疏散 & 基础客流叠加多线路列车释放客流 & 检验多来源客流集中到达时，AA 是否缓解共享设施排队和尾部滞留 \\
\bottomrule
\end{tabular}
\end{table}

\subsection{5.2 网络流分配与瓶颈变化}

实验首先分析网络流分配结果。重点包括各出口通过量、关键共享设施通过量、来源组出口分布和换乘相关客流的跨线流向。如果 AA 有效，其结果不应只表现为总疏散时间变化，还应表现为客流从少数未来排队严重的设施转移到其他可用通道或出口，从而改变关键瓶颈的加载方式。

\begin{table}[htbp]
\centering
\caption{网络流分配与瓶颈变化记录表}
\small
\begin{tabular}{p{0.18\textwidth}p{0.16\textwidth}p{0.20\textwidth}p{0.20\textwidth}p{0.14\textwidth}}
\toprule
场景 & 方法 & 出口利用变化 & 关键设施负荷变化 & 解释 \\
\midrule
低负荷 & 当前状态 A* & 待冻结 & 待冻结 & 基准 \\
低负荷 & AA & 待冻结 & 待冻结 & 稳定性 \\
高负荷 & 当前状态 A* & 待冻结 & 待冻结 & 瓶颈识别 \\
高负荷 & AA & 待冻结 & 待冻结 & 分流效果 \\
\bottomrule
\end{tabular}
\end{table}

该表的解释应服务于后续效率分析。若高负荷场景中 AA 提高出口或关键设施负荷均衡性，并降低少数设施的队列峰值，则可以说明队列预测改变了网络流分配；若分配变化不明显，则即使总疏散时间有差异，也不能轻易归因于 AA 的队列感知机制。

\subsection{5.3 全站疏散效率与方法代价}

在确认网络流分配变化后，再比较整体疏散效率。核心指标包括 \(T_{95}\)、\(T_{99}\)、\(T_{100}\)、累计静止人秒、平均疏散时间、出口 Jain 指数、关键设施 Jain 指数、总移动距离和计算耗时。低负荷场景主要检验方法稳定性；高负荷场景主要检验 AA 是否缩短尾部清空并减少等待暴露。

\begin{table}[htbp]
\centering
\caption{全站疏散效率与方法代价}
\small
\begin{tabular}{p{0.14\textwidth}p{0.16\textwidth}p{0.11\textwidth}p{0.11\textwidth}p{0.11\textwidth}p{0.15\textwidth}p{0.14\textwidth}}
\toprule
场景 & 方法 & \(T_{95}\) & \(T_{99}\) & \(T_{100}\) & 累计静止人秒 & 方法代价 \\
\midrule
低负荷 & 当前状态 A* & 待冻结 & 待冻结 & 待冻结 & 待冻结 & 待冻结 \\
低负荷 & AA & 待冻结 & 待冻结 & 待冻结 & 待冻结 & 待冻结 \\
高负荷 & 当前状态 A* & 待冻结 & 待冻结 & 待冻结 & 待冻结 & 待冻结 \\
高负荷 & AA & 待冻结 & 待冻结 & 待冻结 & 待冻结 & 待冻结 \\
\bottomrule
\end{tabular}
\end{table}

结果冻结后，分析应按“效率--瓶颈--代价”的顺序展开。若 AA 缩短高负荷场景的尾部清空时间并降低累计静止人秒，同时增加总移动距离或计算耗时，应表述为“通过绕开未来队列瓶颈换取等待暴露降低”，而不是写成所有指标全面占优。若低负荷场景中两种方法总清空时间接近，则说明在瓶颈较弱时未来队列预测的边际收益有限，这是方法适用边界的一部分。

\subsection{5.4 Pathfinder 对照与敏感性分析}

Pathfinder 微观仿真用于检查图模型结论在空间层面的趋势表现。对照重点包括累计疏散曲线、拥堵位置、出口利用和来源组滞留特征。由于 Pathfinder 的个体避让、局部冲突和空间占用机制与网络流模型不同，本文不把 Pathfinder 写成逐个乘客轨迹校准，而将其作为微观趋势对照。若两类模型都显示相同关键设施出现拥堵，且疏散曲线尾部变化方向一致，则可认为 Pathfinder 支持网络流模型的主要结论。

敏感性分析应围绕本文机制设置，而不是泛泛调整速度--密度参数。建议重点考察三类因素：换乘客流比例、列车到达重叠程度和关键共享设施服务能力。换乘客流比例用于检验跨线客流增强时共享设施竞争是否更突出；列车到达重叠程度用于检验多来源客流集中释放时 AA 是否更有价值；关键设施服务能力用于检验瓶颈强弱变化下方法的适用范围。

\begin{table}[htbp]
\centering
\caption{敏感性分析设置}
\small
\begin{tabular}{p{0.24\textwidth}p{0.32\textwidth}p{0.32\textwidth}}
\toprule
因素 & 调整方式 & 解释目的 \\
\midrule
换乘客流比例 & 调整换乘来源组占比 & 判断换乘流增强时 AA 是否更能缓解共享瓶颈 \\
列车到达重叠程度 & 调整多线路列车释放时间间隔 & 判断到达高峰越集中时，队列预测是否更重要 \\
关键共享设施服务能力 & 调整楼扶梯、通道、闸机或出口前区容量 & 判断瓶颈强弱变化下方法何时有效、何时收益有限 \\
\bottomrule
\end{tabular}
\end{table}

通过上述安排，第 5 章不再只是比较 AA 与当前状态 A* 的结果表，而是形成完整验证链条：先证明高负荷下网络流在共享设施处形成瓶颈，再证明 AA 改变了到达负荷分配，随后证明这种分配变化改善全站疏散效率，最后说明方法的代价和适用边界。

\end{document}
